% unit: u29
% sources: sources/073-symmetric-matrices-and-positive-definiteness.txt, sources/049-lecture-25-symmetric-matrices-and-positive-definiteness.txt, sources/206-mit18-06scf11-ses3-1sum-pdf-download-file.txt, sources/050-problem-solving-symmetric-matrices-and-positive-definiteness.txt, sources/207-mit18-06scf11-ses3-1prob-pdf-download-file.txt, sources/208-mit18-06scf11-ses3-1sol-pdf-download-file.txt
\section{Symmetric matrices and positive definiteness}\label{sec:u29}

Part~II ended with the eigenvalue picture of a square matrix: when $A$ has $n$ independent
eigenvectors, $A=X\Lambda X^{-1}$ (Section~\ref{sec:u24}). A matrix with special structure should
have special eigenvalues and eigenvectors; Markov matrices, for instance, always have
$\lambda=1$ (Section~\ref{sec:u26}). This unit studies the most important special class,
the \emph{symmetric} matrices $A=A\T$. Two facts carry the whole unit:
\begin{itemize}
  \item the eigenvalues of a real symmetric matrix are \textbf{real};
  \item its eigenvectors can be chosen \textbf{perpendicular}, so $X$ becomes an orthogonal matrix $Q$.
\end{itemize}
Once eigenvalues are known to be real we can ask for their \emph{signs}, and a surprising
shortcut appears: the signs of the pivots are the signs of the eigenvalues. The symmetric
matrices whose eigenvalues are all positive, the positive definite matrices, are where pivots,
determinants and eigenvalues finally come together.

\subsection{Real eigenvalues}

\begin{definition}[Symmetric matrix]\label{def:u29-symmetric}
A square real matrix $A$ is \emph{symmetric} if $A\T=A$, that is $a_{ij}=a_{ji}$ for all $i,j$.
\end{definition}

A real matrix can have complex eigenvalues: the $90^\circ$ rotation
$\left[\begin{smallmatrix}0&-1\\1&0\end{smallmatrix}\right]$ has $\lambda=\pm i$. So to prove that
symmetric matrices do not, we must allow $\lambda$ and $\vb{x}$ to be complex for a moment
and see what symmetry forces. The tool is the complex conjugate: if $z=a+ib$ then
$\bar z=a-ib$ and $\bar z z=a^2+b^2\ge 0$.

\begin{lemma}[Conjugate length]\label{lem:u29-xbarx}
For a complex vector $\vb{x}=(x_1,\dots,x_n)$,
\[
  \bar{\vb{x}}\T\vb{x}=\bar x_1x_1+\cdots+\bar x_nx_n=|x_1|^2+\cdots+|x_n|^2 ,
\]
which is real, and positive unless $\vb{x}=\vb{0}$.
\end{lemma}
\begin{proof}
Each term is $\bar x_kx_k=(a_k-ib_k)(a_k+ib_k)=a_k^2+b_k^2=|x_k|^2$: the cross terms $\pm ia_kb_k$
cancel. A sum of such squares is zero only when every $a_k$ and $b_k$ is zero.
\end{proof}

Without the bar the lemma fails: for $\vb{x}=(1,i)$ we get $\vb{x}\T\vb{x}=1+i^2=0$ although
$\vb{x}\ne\vb{0}$. Section~\ref{sec:u30} builds the geometry of complex vectors on this
observation.

\begin{theorem}[Real eigenvalues]\label{thm:u29-real}
Every eigenvalue of a real symmetric matrix is real.
\end{theorem}
\begin{proof}
Let $A\vb{x}=\lambda\vb{x}$ with $\vb{x}\ne\vb{0}$, where $\lambda$ and $\vb{x}$ may be complex.
Take complex conjugates of every factor. Because $A$ is real, $\bar A=A$ and
\[
  A\bar{\vb{x}}=\bar\lambda\,\bar{\vb{x}}.
\]
(This step alone shows that the complex eigenvalues of any real matrix come in conjugate pairs
$\lambda,\bar\lambda$.) Now transpose and use $A\T=A$, which is the only place symmetry enters:
\[
  \bar{\vb{x}}\T A=\bar\lambda\,\bar{\vb{x}}\T .
\]
Multiply this on the right by $\vb{x}$, and multiply the original equation $A\vb{x}=\lambda\vb{x}$
on the left by $\bar{\vb{x}}\T$:
\[
  \bar{\vb{x}}\T A\vb{x}=\bar\lambda\,\bar{\vb{x}}\T\vb{x},
  \qquad
  \bar{\vb{x}}\T A\vb{x}=\lambda\,\bar{\vb{x}}\T\vb{x}.
\]
The left sides agree, so $(\lambda-\bar\lambda)\,\bar{\vb{x}}\T\vb{x}=0$. By
Lemma~\ref{lem:u29-xbarx}, $\bar{\vb{x}}\T\vb{x}>0$, so $\lambda=\bar\lambda$: the imaginary part
of $\lambda$ is zero.
\end{proof}

Once $\lambda$ is real, $A-\lambda I$ is a real singular matrix, so its nullspace contains real
nonzero vectors. Eigenvectors of a real symmetric matrix can therefore always be taken real.

\begin{remark}
If $A$ is complex, the same proof goes through provided $\bar A\T=A$: transposing
$\bar A\bar{\vb{x}}=\bar\lambda\bar{\vb{x}}$ gives $\bar{\vb{x}}\T\bar A\T=\bar\lambda\bar{\vb{x}}\T$,
and we need $\bar A\T=A$ to continue. So the right complex analogue of ``symmetric'' is
``equal to its conjugate transpose''. These Hermitian matrices are the subject of
Section~\ref{sec:u30}.
\end{remark}

\subsection{Perpendicular eigenvectors and the spectral theorem}

\begin{theorem}[Orthogonal eigenvectors]\label{thm:u29-orth}
If $A$ is real symmetric, $A\vb{x}=\lambda_1\vb{x}$ and $A\vb{y}=\lambda_2\vb{y}$ with
$\lambda_1\ne\lambda_2$, then $\vb{x}\T\vb{y}=0$.
\end{theorem}
\begin{proof}
Compute $\vb{x}\T A\vb{y}$ in two ways. Using $A\vb{y}=\lambda_2\vb{y}$ gives $\lambda_2\vb{x}\T\vb{y}$.
Using symmetry, $\vb{x}\T A\vb{y}=(A\T\vb{x})\T\vb{y}=(A\vb{x})\T\vb{y}=\lambda_1\vb{x}\T\vb{y}$.
Hence $(\lambda_1-\lambda_2)\vb{x}\T\vb{y}=0$, and since $\lambda_1\ne\lambda_2$, $\vb{x}\T\vb{y}=0$.
\end{proof}

When an eigenvalue is repeated there is a whole plane (or larger subspace) of eigenvectors,
and not every pair in it is perpendicular. The identity matrix is the extreme case: every
vector is an eigenvector. The correct statement is that the eigenvectors \emph{can be chosen}
perpendicular, and that there are always enough of them. That is the content of the main
theorem.

\begin{theorem}[Spectral theorem; principal axis theorem]\label{thm:u29-spectral}
Every real symmetric $n\times n$ matrix can be factored as
\[
  A=Q\Lambda Q\T ,
\]
where $Q$ is an orthogonal matrix ($Q\T Q=I$) whose columns are orthonormal eigenvectors
$\vb{q}_1,\dots,\vb{q}_n$ of $A$, and $\Lambda$ is the real diagonal matrix of eigenvalues.
Conversely, every matrix of the form $Q\Lambda Q\T$ with $\Lambda$ real diagonal is symmetric.
\end{theorem}
\begin{proof}
The converse is one line: $(Q\Lambda Q\T)\T=(Q\T)\T\Lambda\T Q\T=Q\Lambda Q\T$.

For the factorization we use induction on $n$; the case $n=1$ is trivial. The characteristic
polynomial of $A$ has at least one root $\lambda_1$, which is real by
Theorem~\ref{thm:u29-real}, with a real unit eigenvector $\vb{q}_1$. Extend $\vb{q}_1$ to an
orthonormal basis $\vb{q}_1,\vb{p}_2,\dots,\vb{p}_n$ of $\R^n$ (Gram--Schmidt,
Section~\ref{sec:u19}) and let $Q_1$ be the orthogonal matrix with these columns. The first
column of $Q_1\T AQ_1$ is $Q_1\T A\vb{q}_1=\lambda_1Q_1\T\vb{q}_1=\lambda_1\vb{e}_1$. The matrix
$Q_1\T AQ_1$ is symmetric, because $(Q_1\T AQ_1)\T=Q_1\T A\T Q_1=Q_1\T AQ_1$, so its first row is
also $\lambda_1\vb{e}_1\T$:
\[
  Q_1\T AQ_1=\begin{bmatrix}\lambda_1&\vb{0}\T\\ \vb{0}&A_2\end{bmatrix},
  \qquad A_2 \text{ symmetric of size } n-1 .
\]
By induction $A_2=Q_2\Lambda_2Q_2\T$ with $Q_2$ orthogonal. Put
$Q=Q_1\left[\begin{smallmatrix}1&\vb{0}\T\\ \vb{0}&Q_2\end{smallmatrix}\right]$, a product of
orthogonal matrices and hence orthogonal. Then
$Q\T AQ=\left[\begin{smallmatrix}\lambda_1&\vb{0}\T\\ \vb{0}&\Lambda_2\end{smallmatrix}\right]=\Lambda$,
that is $A=Q\Lambda Q\T$. Reading $AQ=Q\Lambda$ column by column shows $A\vb{q}_k=\lambda_k\vb{q}_k$.
\end{proof}

Compared with the general diagonalization $A=X\Lambda X^{-1}$, two things have improved:
the eigenvalues in $\Lambda$ are real, and the eigenvector matrix $X$ has become an orthogonal
$Q$, whose inverse is simply $Q\T$. Mathematicians call the set of eigenvalues the
\emph{spectrum}, hence the name; in mechanics the same statement is the principal axis
theorem: in the right axes (the eigenvectors) a symmetric operation does not couple the
directions together.

\begin{corollary}[A sum of projections]\label{cor:u29-projections}
If $A=Q\Lambda Q\T$ then, multiplying columns times rows,
\[
  A=\lambda_1\vb{q}_1\vb{q}_1\T+\lambda_2\vb{q}_2\vb{q}_2\T+\cdots+\lambda_n\vb{q}_n\vb{q}_n\T .
\]
Each $P_k=\vb{q}_k\vb{q}_k\T$ is the projection matrix onto the line through $\vb{q}_k$; these
projections satisfy $P_k^2=P_k$, $P_jP_k=0$ for $j\ne k$, and $P_1+\cdots+P_n=QQ\T=I$.
\end{corollary}
\begin{proof}
The product $Q(\Lambda Q\T)$ is a sum of (column $k$ of $Q$) times (row $k$ of $\Lambda Q\T$),
and row $k$ of $\Lambda Q\T$ is $\lambda_k\vb{q}_k\T$. For a unit vector the projection formula
$\vb{a}\vb{a}\T/\vb{a}\T\vb{a}$ of Section~\ref{sec:u17} reduces to $\vb{q}_k\vb{q}_k\T$. Finally
$P_jP_k=\vb{q}_j(\vb{q}_j\T\vb{q}_k)\vb{q}_k\T$ is $P_k$ when $j=k$ and $0$ otherwise.
\end{proof}

So a symmetric matrix acts on $\vb{x}$ in three moves: split $\vb{x}$ into its pieces
$(\vb{q}_k\T\vb{x})\vb{q}_k$ along perpendicular axes, stretch piece $k$ by $\lambda_k$, add.

\begin{example}[A spectral decomposition, worked in full]\label{ex:u29-spectral}
Take $A=\begin{bmatrix}1&2\\2&-2\end{bmatrix}$. The trace is $-1$ and the determinant is
$-2-4=-6$, so $\lambda^2+\lambda-6=(\lambda-2)(\lambda+3)=0$ and $\lambda_1=2$, $\lambda_2=-3$:
real, as promised.

\emph{Eigenvectors.} $A-2I=\left[\begin{smallmatrix}-1&2\\2&-4\end{smallmatrix}\right]$ has nullspace
spanned by $(2,1)$; $A+3I=\left[\begin{smallmatrix}4&2\\2&1\end{smallmatrix}\right]$ has nullspace
spanned by $(1,-2)$. Their dot product is $2-2=0$, as Theorem~\ref{thm:u29-orth} predicts.
Normalizing,
\[
  \vb{q}_1=\tfrac{1}{\sqrt5}\begin{bmatrix}2\\1\end{bmatrix},\quad
  \vb{q}_2=\tfrac{1}{\sqrt5}\begin{bmatrix}1\\-2\end{bmatrix},\quad
  A=Q\Lambda Q\T=\tfrac{1}{\sqrt5}\begin{bmatrix}2&1\\1&-2\end{bmatrix}
  \begin{bmatrix}2&0\\0&-3\end{bmatrix}
  \tfrac{1}{\sqrt5}\begin{bmatrix}2&1\\1&-2\end{bmatrix}.
\]
\emph{Projections.}
\[
  \vb{q}_1\vb{q}_1\T=\tfrac15\begin{bmatrix}4&2\\2&1\end{bmatrix},\qquad
  \vb{q}_2\vb{q}_2\T=\tfrac15\begin{bmatrix}1&-2\\-2&4\end{bmatrix},\qquad
  2\cdot\tfrac15\begin{bmatrix}4&2\\2&1\end{bmatrix}-3\cdot\tfrac15\begin{bmatrix}1&-2\\-2&4\end{bmatrix}
  =\tfrac15\begin{bmatrix}5&10\\10&-10\end{bmatrix}=A .
\]
\emph{Action on a vector.} For $\vb{x}=(1,1)$ the two pieces are
$\vb{q}_1\vb{q}_1\T\vb{x}=(1.2,\,0.6)$ and $\vb{q}_2\vb{q}_2\T\vb{x}=(-0.2,\,0.4)$, which add back to
$\vb{x}$. Stretching them by $2$ and $-3$ gives $(2.4,1.2)+(0.6,-1.2)=(3,0)$, and indeed
$A\vb{x}=(1+2,\,2-2)=(3,0)$. Figure~\ref{fig:u29-projections} draws this.

\emph{Pivots.} Elimination subtracts $2$ times row~1 from row~2, leaving pivots $1$ and
$-2-4=-6$: one positive and one negative, exactly like the eigenvalues $2$ and $-3$. Their
products agree too: $1\cdot(-6)=2\cdot(-3)=\det A$.
\end{example}

\begin{figure}[ht]\centering
\begin{tikzpicture}[scale=1.55,>=Stealth]
  \draw[gray!60,dashed] (-0.9,-0.45) -- (3.4,1.7) node[above,black] {line of $\vb{q}_1$};
  \draw[gray!60,dashed] (0.35,-0.7) -- (-0.7,1.4) node[above,black] {line of $\vb{q}_2$};
  \draw[->,gray] (-0.9,0) -- (3.6,0);
  \draw[->,gray] (0,-0.7) -- (0,1.6);
  \draw[dotted,thick] (1,1) -- (1.2,0.6);
  \draw[dotted,thick] (1,1) -- (-0.2,0.4);
  \draw[->,very thick] (0,0) -- (1,1) node[above] {$\vb{x}$};
  \draw[->,thick,blue] (0,0) -- (1.2,0.6) node[below right=-1pt] {$P_1\vb{x}$};
  \draw[->,thick,orange!85!black] (0,0) -- (-0.2,0.4) node[left] {$P_2\vb{x}$};
  \draw[->,thick,blue!60] (1.2,0.6) -- (2.4,1.2) node[above left] {$2P_1\vb{x}$};
  \draw[->,thick,orange!70!black] (2.4,1.2) -- (3,0);
  \node[orange!70!black,right] at (2.75,0.7) {$-3P_2\vb{x}$};
  \draw[->,very thick,red!75!black] (0,0) -- (3,0) node[below] {$A\vb{x}=(3,0)$};
  \draw (0.112,0.056) -- (0.068,0.144) -- (-0.044,0.088);
\end{tikzpicture}
\caption{The spectral theorem at work for $A=\left[\begin{smallmatrix}1&2\\2&-2\end{smallmatrix}\right]$
and $\vb{x}=(1,1)$. The eigenvector lines are perpendicular; $\vb{x}$ splits into its projections
$P_1\vb{x}=(1.2,0.6)$ and $P_2\vb{x}=(-0.2,0.4)$; $A$ stretches them by $\lambda_1=2$ and
$\lambda_2=-3$ and adds, giving $A\vb{x}=(3,0)$.}
\label{fig:u29-projections}
\end{figure}

\subsection{Signs of pivots are signs of eigenvalues}

Knowing the eigenvalues are real, the next question is whether they are positive or negative;
for $d\vb{u}/dt=A\vb{u}$ that decides stability (Section~\ref{sec:u25}). Finding the 50
eigenvalues of a $50\times 50$ matrix from $\det(A-\lambda I)=0$ is hopeless: the roots of a
degree-50 polynomial are extremely sensitive to rounding. Pivots, on the other hand, are cheap
and safe to compute. For symmetric matrices they carry the sign information we want.

If elimination on a symmetric $A$ needs no row exchanges, it gives $A=LDU$ with $L$ unit lower
triangular, $D$ the diagonal matrix of pivots and $U$ unit upper triangular
(Section~\ref{sec:u05}). Transposing gives a second such factorization $A=A\T=U\T DL\T$, and
since the $LDU$ factors are unique, $U=L\T$:
\[
  A=LDL\T\qquad (A \text{ symmetric, no row exchanges}).
\]

\begin{theorem}[Pivot signs]\label{thm:u29-signs}
Let $A$ be real symmetric with nonzero pivots $d_1,\dots,d_n$ (no row exchanges). Then the
number of positive pivots equals the number of positive eigenvalues, and likewise for
negative ones.
\end{theorem}
\begin{proof}
Write $A=LDL\T$ and, for $0\le t\le 1$, let $L_t=tL+(1-t)I$. Every $L_t$ is lower triangular with
ones on the diagonal, so $\det L_t=1$. The matrix $A_t=L_tDL_t\T$ is symmetric, so its eigenvalues
are real, and $\det A_t=\det D=d_1\cdots d_n\ne 0$, so $0$ is never an eigenvalue of $A_t$. At
$t=0$, $A_0=D$, whose eigenvalues are the pivots; at $t=1$, $A_1=A$. The eigenvalues of $A_t$ are
the roots of a polynomial whose coefficients depend continuously on $t$, so they move
continuously. A real eigenvalue that changed sign would have to pass through $0$, which never
happens. Hence $A$ and $D$ have the same number of positive eigenvalues.
\end{proof}

The product of the pivots and the product of the eigenvalues are both $\det A$, but that is
one number; the theorem gives $n$ separate pieces of sign information. Combined with a shift
it locates eigenvalues: the eigenvalues of $A-bI$ are $\lambda_k-b$, so the positive pivots of
$A-bI$ count the eigenvalues of $A$ that are larger than $b$.

\begin{example}[Counting eigenvalues above a number]\label{ex:u29-shift}
Let $K=\begin{bmatrix}2&-1&0\\-1&2&-1\\0&-1&2\end{bmatrix}$. How many eigenvalues exceed $5/2$?
Eliminate on
\[
  K-\tfrac52 I=\begin{bmatrix}-\tfrac12&-1&0\\-1&-\tfrac12&-1\\0&-1&-\tfrac12\end{bmatrix}.
\]
The first pivot is $-\tfrac12$. Subtracting $2\times$row~1 from row~2 leaves
$(0,\ \tfrac32,\ -1)$, so the second pivot is $\tfrac32$. Adding $\tfrac23\times$ the new row~2
to row~3 leaves $-\tfrac12-\tfrac23=-\tfrac76$. The pivots $-\tfrac12,\ \tfrac32,\ -\tfrac76$ have one
positive sign, so exactly one eigenvalue of $K$ is above $5/2$. Check: the eigenvalues of $K$
are $2-\sqrt2\approx0.59$, $2$ and $2+\sqrt2\approx3.41$ (Section~\ref{sec:u31}), and only
$2+\sqrt2$ exceeds $2.5$.
\end{example}

\begin{trap}
Pivots are not eigenvalues. For $A=\left[\begin{smallmatrix}5&2\\2&3\end{smallmatrix}\right]$ the
pivots are $5$ and $11/5$, the eigenvalues are $4\pm\sqrt5\approx 6.24,\,1.76$. What they share is
their signs and their product ($\det A=11$). Even the signs agree only for symmetric
matrices: the nonsymmetric $\left[\begin{smallmatrix}1&-6\\1&-4\end{smallmatrix}\right]$ has
positive pivots $1$ and $2$, but its eigenvalues are $-1$ and $-2$.
\end{trap}

\subsection{Positive definite matrices}

\begin{definition}[Positive definite]\label{def:u29-pd}
A \emph{positive definite} matrix is a real symmetric matrix whose eigenvalues are all
positive.
\end{definition}

Symmetric matrices are good; positive definite ones are the best. Testing the definition
directly needs the eigenvalues, but pivots and determinants do the job without them.

\begin{theorem}[Three equivalent tests]\label{thm:u29-pdtests}
For a real symmetric $n\times n$ matrix $A$ the following are equivalent:
\begin{enumerate}[(a)]
  \item all eigenvalues of $A$ are positive;
  \item all $n$ upper-left determinants $D_1=a_{11},\ D_2,\ \dots,\ D_n=\det A$ are positive,
        where $D_k$ is the determinant of the top-left $k\times k$ submatrix $A_k$;
  \item elimination without row exchanges produces $n$ pivots, all positive.
\end{enumerate}
\end{theorem}
\begin{proof}
(a)$\Rightarrow$(b). Write $A=Q\Lambda Q\T$ and $\vb{y}=Q\T\vb{x}$. Then
$\vb{x}\T A\vb{x}=\vb{y}\T\Lambda\vb{y}=\lambda_1y_1^2+\cdots+\lambda_ny_n^2$, which is positive for every
$\vb{x}\ne\vb{0}$ (since then $\vb{y}\ne\vb{0}$). Now take $\vb{z}\in\R^k$, $\vb{z}\ne\vb{0}$, and pad it with
zeros to $\vb{x}=(\vb{z},0,\dots,0)$; then $\vb{z}\T A_k\vb{z}=\vb{x}\T A\vb{x}>0$. If $\mu$ is an eigenvalue of
the symmetric matrix $A_k$ with real eigenvector $\vb{z}$, this says $\mu\,\vb{z}\T\vb{z}>0$, so $\mu>0$.
Hence $D_k=\det A_k$, the product of the eigenvalues of $A_k$, is positive.

(b)$\Rightarrow$(c). Elimination on the first $k$ rows only ever uses entries of $A_k$, so the
first $k$ pivots of $A$ are the pivots of $A_k$ and $D_k=d_1d_2\cdots d_k$. If every $D_k>0$, then
$d_1=D_1>0$ and $d_k=D_k/D_{k-1}>0$; in particular no zero pivot forces a row exchange.

(c)$\Rightarrow$(a). With all pivots positive, Theorem~\ref{thm:u29-signs} says all $n$
eigenvalues are positive.
\end{proof}

The quantity $\vb{x}\T A\vb{x}$ that appeared in the first step will become the central character
of Section~\ref{sec:u31}, where it gives a fourth test and the link to minima.

\begin{example}[The lecture's example]\label{ex:u29-532}
$A=\begin{bmatrix}5&2\\2&3\end{bmatrix}$ is symmetric. Its upper-left determinants are $5$ and
$15-4=11$, both positive. Its pivots are $5$ and $11/5$ (the second pivot is $\det A$ divided by
the first). So $A$ is positive definite, and the eigenvalues must be positive without our
computing them. Computing anyway: $\lambda^2-8\lambda+11=0$ gives $\lambda=4\pm\sqrt{16-11}=4\pm\sqrt5$,
and $4-\sqrt5\approx1.76>0$.

By contrast $B=\begin{bmatrix}-1&0\\0&-3\end{bmatrix}$ has $\det B=3>0$, yet its eigenvalues
(and pivots) are $-1$ and $-3$. It fails test~(b) at the first step: $D_1=-1<0$.
\end{example}

\begin{trap}
A positive determinant does not make a symmetric matrix positive definite: two negative
eigenvalues also multiply to a positive number. Test \emph{all} $n$ upper-left determinants,
not only the last one. Another quick necessary condition: every diagonal entry
$a_{ii}=\vb{e}_i\T A\vb{e}_i$ of a positive definite matrix is positive, so a single
$a_{ii}\le0$ rules it out at once.
\end{trap}

Some immediate consequences, each read off from the eigenvalues:
\begin{itemize}
  \item a positive definite matrix is invertible, since $\det A=\lambda_1\cdots\lambda_n>0$;
  \item a diagonal matrix with positive diagonal entries is positive definite, since those entries
        are its eigenvalues;
  \item the inverse of a positive definite matrix is positive definite, with eigenvalues
        $1/\lambda_k>0$ (more of this in Section~\ref{sec:u32}).
\end{itemize}

\subsection{Exercises}

\begin{exercise}[Where a false proof breaks]\label{exr:u29-false}
Here is a ``proof'' that every real matrix has real eigenvalues: from $A\vb{x}=\lambda\vb{x}$ we get
$\vb{x}\T A\vb{x}=\lambda\,\vb{x}\T\vb{x}$, so $\lambda=\vb{x}\T A\vb{x}/\vb{x}\T\vb{x}$, a ratio of real numbers.
Test each step on the rotation $A=\left[\begin{smallmatrix}0&-1\\1&0\end{smallmatrix}\right]$ with
$\lambda=i$ and $\vb{x}=(i,1)$, and find the hidden assumption.
\end{exercise}
\begin{solution}
First $A\vb{x}=(0\cdot i-1,\ i+0)=(-1,\,i)=i\,(i,1)$, so $\lambda=i$ really is an eigenvalue. Next,
$\vb{x}\T A\vb{x}=(i,1)\cdot(-1,i)=-i+i=0$ and $\lambda\,\vb{x}\T\vb{x}=i\,(i^2+1)=0$, so the equation
$\vb{x}\T A\vb{x}=\lambda\vb{x}\T\vb{x}$ holds, as $0=0$. The final step divides by $\vb{x}\T\vb{x}=i^2+1^2=0$.
The hidden assumption is $\vb{x}\T\vb{x}\ne0$, which is true for real vectors but false for complex
ones; and the ratio is not obviously real either, since $\vb{x}$ may be complex. The correct
version uses $\bar{\vb{x}}\T\vb{x}=|x_1|^2+|x_2|^2>0$ (Lemma~\ref{lem:u29-xbarx}) and needs symmetry.
\end{solution}

\begin{exercise}[Which families are groups?]\label{exr:u29-groups}
A family of invertible matrices is a \emph{group} if it contains $AB$ and $A^{-1}$ whenever it
contains $A$ and $B$. Which of these are groups?
\begin{enumerate}[(a)]
  \item symmetric positive definite matrices;
  \item orthogonal matrices;
  \item all exponentials $e^{tA}$ of one fixed matrix $A$, $t\in\R$;
  \item matrices with determinant~1.
\end{enumerate}
\end{exercise}
\begin{solution}
\begin{enumerate}[(a)]
  \item No. $A=\left[\begin{smallmatrix}2&1\\1&1\end{smallmatrix}\right]$ (upper-left determinants $2,1$) and
        $B=\left[\begin{smallmatrix}1&1/2\\1/2&1\end{smallmatrix}\right]$ (determinants $1,3/4$) are positive
        definite, but $AB=\left[\begin{smallmatrix}2.5&2\\1.5&1.5\end{smallmatrix}\right]$ is not even
        symmetric. (Inverses do stay positive definite.)
  \item Yes. If $A\T A=I$ then $A^{-1}=A\T$ and $(A\T)\T A\T=AA\T=I$; if also $B\T B=I$ then
        $(AB)\T(AB)=B\T A\T AB=B\T B=I$.
  \item Yes. $(e^{pA})^{-1}=e^{-pA}$ and $e^{pA}e^{qA}=e^{(p+q)A}$ are again of the form $e^{tA}$.
  \item Yes. $\det(A^{-1})=1/\det A=1$ and $\det(AB)=\det A\det B=1$.
\end{enumerate}
\end{solution}

\begin{exercise}[Projections and determinants]\label{exr:u29-projection}
\begin{enumerate}[(a)]
  \item Show that the only positive definite projection matrix is $P=I$.
  \item Show that $S=\left[\begin{smallmatrix}-3&1\\1&-2\end{smallmatrix}\right]$ has $\det S>0$ but is not
        positive definite, by exhibiting $\vb{x}$ with $\vb{x}\T S\vb{x}<0$.
\end{enumerate}
\end{exercise}
\begin{solution}
\begin{enumerate}[(a)]
  \item A projection satisfies $P^2=P$, so each eigenvalue satisfies $\lambda^2=\lambda$: $\lambda\in\{0,1\}$.
        Positive definite forces $\lambda>0$, so all eigenvalues are $1$. A projection matrix is
        symmetric, so $P=Q\Lambda Q\T=QIQ\T=I$.
  \item $\det S=6-1=5>0$, but $\vb{x}=(1,0)$ gives $\vb{x}\T S\vb{x}=s_{11}=-3<0$. In the language of
        Theorem~\ref{thm:u29-pdtests}, $D_1=-3$ fails. In fact both eigenvalues,
        $(-5\pm\sqrt5)/2\approx-1.38,\,-3.62$, are negative: $S$ is negative definite.
\end{enumerate}
\end{solution}

\begin{exercise}[Spectral decomposition and pivots]\label{exr:u29-345}
For $A=\left[\begin{smallmatrix}3&4\\4&-3\end{smallmatrix}\right]$ find $\lambda_1,\lambda_2$, write
$A=\lambda_1P_1+\lambda_2P_2$ with projections $P_k=\vb{q}_k\vb{q}_k\T$, and confirm that the pivot signs
match the eigenvalue signs.
\end{exercise}
\begin{solution}
Trace $0$ and determinant $-9-16=-25$ give $\lambda^2-25=0$, so $\lambda=\pm5$. From
$A-5I=\left[\begin{smallmatrix}-2&4\\4&-8\end{smallmatrix}\right]$ the eigenvector is $(2,1)$; from
$A+5I=\left[\begin{smallmatrix}8&4\\4&2\end{smallmatrix}\right]$ it is $(1,-2)$; they are perpendicular.
With $\vb{q}_1=(2,1)/\sqrt5$, $\vb{q}_2=(1,-2)/\sqrt5$,
\[
  5\cdot\tfrac15\begin{bmatrix}4&2\\2&1\end{bmatrix}-5\cdot\tfrac15\begin{bmatrix}1&-2\\-2&4\end{bmatrix}
  =\begin{bmatrix}3&4\\4&-3\end{bmatrix}.
\]
Elimination subtracts $\tfrac43$ of row~1 from row~2, so the pivots are $3$ and
$-3-\tfrac{16}{3}=-\tfrac{25}{3}$: one positive, one negative, like $5$ and $-5$. Their product is
$-25=\det A$.
\end{solution}
